\documentclass[11pt,a4paper]{article}
\usepackage[T1]{fontenc}
\usepackage[utf8]{inputenc}
\usepackage[french]{babel}
\usepackage{amsmath,amssymb,amsthm}
\usepackage{geometry}
\geometry{margin=2.5cm}
\setlength{\parindent}{0pt}
\newcommand{\normal}{\mathrel{\triangleleft}}
\newcommand{\semi}{\rtimes}

\title{Proposition de corrigé TD1 CM}
\author{Maxence Caucheteux}
\date{October 30, 2023}

\begin{document}
\maketitle

\subsection*{Exercice 2}
\textbf{Q1}

$SL_n(k)$ est un sous-groupe de $GL_n(k)$ (vérifications élémentaires). Pour $A\in GL_n(k)$ et $B\in SL_n(k)$, $\det(ABA^{-1})=\det(B)=1$ donc $ABA^{-1}\in SL_n(k)$. Donc $SL_n(k)\normal GL_n(k)$. L'application
\[
\begin{array}{rcl}
\phi:GL_n(k)/SL_n(k)&\longrightarrow&k^*\\
\overline M&\longmapsto&\det(M)
\end{array}
\]
est bien définie et est un isomorphisme de groupes.

\textbf{Preuve :} Il est clair qu'elle est bien définie (image indépendante du choix du représentant et bien à valeur dans $k^*$) et que c'est un morphisme de groupes.

Si $\phi(\overline M)=\det(M)=1$, alors $M\in SL_n(k)$ et $\overline M=SL_n(k)$, donc $\phi$ est injective. Pour $x\in k^*$, on pose $M=D(x,1_k,\ldots,1_k)\in GL_n(k)$, de sorte que $\phi(\overline M)=\det(M)=x$, ce qui montre la surjectivité de $\phi$.

Ainsi $GL_n(k)/SL_n(k)\simeq k^*$.

\textbf{Q2} Soit $\tau$ une transposition de $H$ (il en existe une par hypothèse). On note $\tau=(i\ j)$. On a :
\[
  \forall\sigma\in S_n,\quad \sigma\tau\sigma^{-1}=(\sigma(i)\ \sigma(j)).
\]
Cela prouve que $H$ contient toutes les transpositions de $S_n$. Comme les transpositions engendrent $S_n$, on a $H=S_n$.

\subsection*{Exercice 3 - Centre}
\textbf{Q1} Les $f_g$ ($g\in G$) sont des automorphismes.

\textbf{Morphismes :} $\forall g\in G$, $\forall x_1,x_2\in G$,
\[
f_g(x_1x_2)=gx_1x_2g^{-1}=(gx_1g^{-1})(gx_2g^{-1})=f_g(x_1)f_g(x_2).
\]
\textbf{Bijectifs :} Les morphismes $f_g$ admettent pour réciproque
\[
\begin{array}{rcl}
\varphi:G&\longrightarrow&G\\
x&\longmapsto&g^{-1}xg.
\end{array}
\]
Ainsi $\operatorname{Int}(G)\subset\operatorname{Aut}(G)$.

$\operatorname{Int}(G)$ est un sous-groupe de $\operatorname{Aut}(G)$.
\[
\operatorname{id}_G=f_e\quad\text{donc}\quad \operatorname{id}_G\in\operatorname{Int}(G).
\]
Si $g_1,g_2\in G$, pour $x\in G$,
\[
f_{g_1}\circ f_{g_2}(x)=f_{g_1}(g_2xg_2^{-1})=g_1g_2xg_2^{-1}g_1^{-1}=(g_1g_2)x(g_1g_2)^{-1}=f_{g_1g_2}(x).
\]
Donc $f_{g_1}\circ f_{g_2}\in\operatorname{Int}(G)$. Si $g\in G$, $f_g^{-1}=f_{g^{-1}}$ donc $f_g^{-1}\in\operatorname{Int}(G)$.

$Z(G)$ est un sous-groupe normal de $G$.

\textbf{Preuve :} C'est bien un sous-groupe de $G$ (vérifications élémentaires) et il est distingué car pour $g\in G$ et $h\in Z(G)$, pour tout $x\in G$, $ghg^{-1}x=gg^{-1}hx=hx$ et $xghg^{-1}=xhgg^{-1}=xh=hx$. Donc $ghg^{-1}\in Z(G)$, ce qui conclut.

\textbf{Q2} On note $(\overline{g_i})_{i\in I}$ les classes d'équivalence distinctes ($g_i\in G$). On définit
\[
\begin{array}{rcl}
\phi:G/Z(G)&\longrightarrow&\operatorname{Int}(G)\\
\overline{g_i}&\longmapsto&f_{g_i}.
\end{array}
\]
$\phi$ est bien définie. En effet, pour $i\in I$ et $z\in Z(G)$,
\[
\forall x\in G,\quad f_{g_iz}(x)=(g_iz)x(g_iz)^{-1}=g_ixg_i^{-1}
\]
(image indépendante du représentant choisi).

Pour $i,j\in I$,
\[
\phi(\overline{g_i}\,\overline{g_j})=\phi(\overline{g_ig_j})=f_{g_ig_j}=f_{g_i}\circ f_{g_j}=\phi(\overline{g_i})\circ\phi(\overline{g_j}).
\]
Donc $\phi$ est un morphisme. Si $\overline{g_i}$ vérifie $\phi(\overline{g_i})=\operatorname{id}_G$, c'est-à-dire
\[
\forall x\in G,\quad g_ix=xg_i,
\]
alors $g_i\in Z(G)$ et $\overline{g_i}=Z(G)$. Ainsi $\phi$ est injective. Soit $f_g\in\operatorname{Int}(G)$, avec $g\in G$; on a $\phi(\overline g)=f_g$. Donc $\phi$ est surjective.

\textbf{Bilan :} $\phi$ est un isomorphisme, donc
\[
G/Z(G)\simeq\operatorname{Int}(G).
\]

\textbf{Q3} Si $G/Z(G)$ est monogène, $G/Z(G)=\langle gZ(G)\rangle$. Soient $g_1,g_2\in G$. Alors $g_1=g^{k_1}z_1$ et $g_2=g^{k_2}z_2$ avec $z_1,z_2\in Z(G)$. Ainsi
\[
g_1g_2=g^{k_1}z_1g^{k_2}z_2=g^{k_1}g^{k_2}z_1z_2=g^{k_2}g^{k_1}z_2z_1=g^{k_2}z_2g^{k_1}z_1=g_2g_1.
\]
Donc $G$ est abélien.

\subsection*{Exercice 4 - Décomposition directe}
On définit
\[
\begin{array}{rcl}
\phi:H\times K&\longrightarrow&G\\
(h,k)&\longmapsto&hk.
\end{array}
\]
$\phi$ est un morphisme de groupes.

\textbf{Preuve :} On montre d'abord que pour $(h,k)\in H\times K$, $hk=kh$. Comme $H\normal G$, $kh^{-1}k^{-1}\in H$ puis $hkh^{-1}k^{-1}\in H$. De même, comme $K\normal G$, $hkh^{-1}\in K$ puis $hkh^{-1}k^{-1}\in K$. Ainsi $hkh^{-1}k^{-1}\in H\cap K=\{e_G\}$. Donc $hk=kh$.

Ainsi, pour $(h_1,k_1),(h_2,k_2)\in H\times K$,
\[
\phi(h_1h_2,k_1k_2)=h_1h_2k_1k_2=h_1k_1h_2k_2=\phi(h_1,k_1)\phi(h_2,k_2).
\]
$\phi$ est donc un morphisme de groupes. Il est surjectif car $HK=G$. Si $(h,k)\in H\times K$ vérifie $\phi(h,k)=hk=e_G$, alors $h=k^{-1}\in H\cap K=\{e_G\}$, donc $h=k=e_G$. Cela montre l'injectivité de $\phi$.

$\phi$ est donc un isomorphisme : $G\simeq H\times K$.

\subsection*{Exercice 5 - Groupe dérivé}
\textbf{Q1} Pour $g\in G$, $z=[x,y]\in D(G)$ (avec $x,y\in G$), on remarque que
\[
gzg^{-1}=[gxg^{-1},gyg^{-1}]\in D(G).
\]
Pour $z\in D(G)$, $z$ s'écrit sous la forme $z=z_1\cdots z_k$ où les $z_i$ sont des commutateurs. Alors pour tout $g\in G$, par ce qui précède,
\[
gzg^{-1}=gz_1\cdots z_kg^{-1}=(gz_1g^{-1})\cdots(gz_kg^{-1})\in D(G).
\]
On a donc montré que $D(G)\normal G$.

\textbf{Q2}

\textbf{Sens direct :} Si $D(G)\subset H$, alors $[g_1^{-1},g_2^{-1}]\in H$, c'est-à-dire $g_1^{-1}g_2^{-1}g_1g_2\in H$, puis $g_1g_2\in g_2g_1H$. Donc $(g_1g_2)H\subset(g_2g_1)H$. Par symétrie des rôles entre $g_1$ et $g_2$, l'inclusion réciproque est vraie et $(g_1g_2)H=(g_2g_1)H$. Ainsi $G/H$ est abélien.

\textbf{Sens indirect :} Si $G/H$ est abélien. Soient $x,y\in G$. On veut montrer que $[x,y]\in H$. On écrit $x=g_1h_1$ et $y=g_2h_2$ avec $g_1,g_2\in G$ et $h_1,h_2\in H$. Alors
\[
[x,y]=g_1h_1g_2h_2(g_2h_2g_1h_1)^{-1}.
\]
Or, comme $G/H$ est abélien, il existe $\widetilde h,\widetilde{\widetilde h}\in H$ tels que $g_1h_1g_2h_2=g_2g_1h_1\widetilde h$ et $g_2h_2g_1h_1=g_1g_2h_2\widetilde{\widetilde h}$. De sorte que
\[
[x,y]=g_2g_1h_1\widetilde h\widetilde{\widetilde h}^{-1}h_2^{-1}g_2^{-1}g_1^{-1}.
\]
Comme $G/H$ est abélien, il existe $h\in H$ tel que $g_1g_2=g_2g_1h$. Ainsi, par ce qui précède et puisque $H\normal G$,
\[
[x,y]=(g_1g_2)(h^{-1}h_1\widetilde h\widetilde{\widetilde h}^{-1}h_2^{-1})(g_1g_2)^{-1}\in H.
\]
Donc $D(G)\subset H$.

\subsection*{Exercice 6 - Groupe résoluble, Galois}
\textbf{Q2}

\textbf{Sens direct}

Si $G$ est résoluble, il existe $n_0\in\mathbb N^*$ tel que
\[
G_0=\{e\}\normal G_1\normal\cdots\normal G_{n_0}=G
\]
avec $\forall i\in[0,n_0-1]$, $G_{i+1}/G_i$ abélien. Par l'exercice 5, comme $G_{i+1}/G_i$ est abélien,
\[
\forall i\in[0,n_0-1],\quad D(G_{i+1})\subset G_i.
\]
Une récurrence immédiate montre alors que $D^{n_0}(G)\subset\{e\}$. Donc $D^{n_0}(G)=\{e\}$, ce qui conclut.

\textbf{Sens indirect}

Si $(D^n(G))_{n\in\mathbb N}$ stationne à $\{e\}$. Soit $n_0$ un entier tel que $D^{n_0}(G)=\{e\}$. Alors
\[
G_0=\{e\}\normal G_1\normal\cdots\normal G_{n_0}=G
\]
et, en vertu de l'exercice 5, $\forall i\in[0,n_0-1]$, $D^{i+1}(G)/D^i(G)$ est abélien. Donc $G$ est résoluble.

\subsection*{Exercice 7 - Exemples de groupes résolubles}
\textbf{Q2} Comme $rs=sr^{-1}$, des calculs élémentaires montrent que $D(D_n)=\langle r^2\rangle$. Ce dernier groupe étant abélien, $D^2(D_n)=\{e\}$. Par l'exercice 6, Q2, $D_n$ est donc résoluble.

\subsection*{Exercice 8 - Groupe simple}
\textbf{Q1} Tout sous-groupe $H$ d'un groupe abélien $G$ est distingué. Ainsi, les groupes abéliens simples sont les groupes $G$ admettant exactement deux sous-groupes, à savoir $\{e_G\}$ et $G$.

Or les seuls groupes $G$ non triviaux dont les seuls sous-groupes sont $\{e_G\}$ et $G$ sont les groupes cycliques d'ordre $p$ premier (ce sont donc les groupes isomorphes à un $\mathbb Z/p\mathbb Z$ avec $p$ premier).

\textbf{Preuve :} Sens indirect immédiat par théorème de Lagrange. Pour le sens direct, si $G$ est non trivial et n'admet que les sous-groupes triviaux, avec $x\in G$ différent de $e_G$, on a alors $\langle x\rangle=G$. $G$ est donc monogène. Si par l'absurde il était infini, il serait isomorphe à $\mathbb Z$, ce qui n'est pas possible car $\mathbb Z$ admet des sous-groupes stricts. $G$ est donc cyclique.

$G$ est en fait d'ordre premier. En effet, si par l'absurde il ne l'est pas, avec $1<d<|G|$ un diviseur de $|G|$, un lemme connu montre alors qu'il existe un élément $x\in G$ d'ordre $d$. Ainsi le groupe $\langle x\rangle$ est un sous-groupe strict de $G$, ce qui est absurde.

\textbf{Q2} Si $G$ est simple et résoluble, il est abélien. Il est donc isomorphe à un $\mathbb Z/p\mathbb Z$ avec $p$ premier par Q1.

\subsection*{Exercice 10}
On définit
\[
\begin{array}{rcl}
\phi:H&\longrightarrow&\operatorname{Aut}(N)\\
h&\longmapsto&\phi_h:n\mapsto hnh^{-1}.
\end{array}
\]
$\phi$ est bien définie car $N$ est distingué dans $G$ et c'est un morphisme de groupes (vérifications élémentaires).

On définit alors
\[
\begin{array}{rcl}
f:N\semi_{\phi}H&\longrightarrow&G\\
(n,h)&\longmapsto&nh.
\end{array}
\]
Des vérifications élémentaires montrent que $f$ est un morphisme de groupes. C'est en fait un isomorphisme de groupes. L'injectivité est vraie grâce à $N\cap H=\{e_G\}$ et la surjectivité provient du fait que $NH=G$. Ainsi, $N\semi_{\phi}H\simeq G$.

\subsection*{Exercice 11}
\textbf{Q1} On note $C=\{(x,y)\in G^2\mid xy=yx\}$. La probabilité recherchée est $\frac{|C|}{|G\times G|}=\frac{|C|}{|G|^2}$. On a
\begin{align}
C&=\coprod_{x\in G}\{x\}\times C_x,\\
 &=\left(\coprod_{x\in Z(G)}\{x\}\times C_x\right)\coprod
 \left(\coprod_{x\notin Z(G)}\{x\}\times C_x\right),\\
 &=\left(\coprod_{x\in Z(G)}\{x\}\times G\right)\coprod
 \left(\coprod_{x\notin Z(G)}\{x\}\times C_x\right).
\end{align}
Puis, en cardinaux,
\begin{align}
|C|&=|Z(G)|\,|G|+\sum_{x\notin Z(G)}|C_x|,\\
 &=zn+\sum_{x\notin Z(G)}|C_x|,\\
 &\leq zn+(|G|-|Z(G)|)\max_{x\notin Z(G)}|C_x|.
\end{align}
Or, les $C_x$ sont des sous-groupes de $G$ (vérifications élémentaires). Pour $x\notin Z(G)$, ce sont des sous-groupes stricts de $G$. Par théorème de Lagrange,
\[
\forall x\notin Z(G),\quad |C_x|\leq\frac n2.
\]
Puis
\[
|C|\leq\frac{zn}{2}+\frac{n^2}{2}.
\]

\textbf{Q2} Comme $G$ n'est pas abélien, via contraposée de l'exercice 3, $G/Z(G)$ n'est pas monogène. On peut montrer que tous les groupes d'ordre $1,2,3$ sont cycliques. Ainsi, $|G/Z(G)|\geq4$, puis $z\leq\frac n4$.

\textbf{Q3} Par suite,
\[
|C|\leq\frac{n^2}{8}+\frac{n^2}{2}=\frac{5n^2}{8}.
\]
Puis
\[
\frac{|C|}{|G|^2}\leq\frac58.
\]

\textbf{Remarque} Si $G$ est abélien, la probabilité recherchée vaut $1$.

\end{document}
